\chapter{采集与资源交接}\label{ch:capture}
\section{书写者与任务管理（0.1.37）}
先选择或新建书写者，新建只填写显示名称，系统自动生成稳定唯一身份；旧 user\_001 原样保留，修改显示名称不改样本归属。汉字与扩展共用作者列表，启动不自动打开任务。导入 TXT/CSV 先集中确认作者、字表、份数与目录；取消不替换当前任务。

\ui{暂停并关闭}保留任务进度；\ui{取消任务}只停止剩余采集，已有样本不删除；在\ui{任务列表}恢复后可继续。列表统一展示汉字、扩展、已完成、已取消及损坏任务，可搜索，支持修改任务名称、归档、复制未录部分到另一个作者。复制不继承原作者的复用确认。零样本任务可经确认更换作者，已有样本不直接改归属。

两类采集均可查看和重录已保存份次，保留原始校验备份，汉字对应旧 NPY/TXT 也从活动数据中退役。已有字体、生成结果和冻结文档不自动改写。未保存笔迹提供保存、放弃或返回；多人窗口发生进度冲突时拒绝覆盖并保留笔迹。Windows 数据目录过深时提示改选较短目录。以下历史操作中填写 ID 的步骤，在当前统一应用中使用书写者选择。
\entry{统一应用左侧\ui{采集}导航，内含两个子页：汉字样本、标点与扩展。}

\section{汉字采集}
填写作者 ID，创建或打开采集任务，使用数位板书写，核对笔迹后保存。汉字单字任务保留标准笔数和笔序检查。画布使用数位笔输入；鼠标可操作按钮和表单，不能代替数位笔录制。应用不再显示整行采集页，已有整行资源仍保留。

完成后切换生成页，填写作者 ID、选择汉字样本目录，依次导入、建档、生成。导入会检查作者身份，不能仅修改界面的作者名来混用其他作者样本。

\section{标点与扩展}
目标字符使用大字预览，正体与斜体分别显示。误录且已经保存时，点击\ui{重录已保存…}，查看旧笔迹并选择对应份次；可以批量选择同字符、同写法的全部已保存记录。确认后所选份次退回待采，原始记录保留备份，正确样本不变。重新保存产生新样本，不会把旧错误记录再次加入字体。\ui{清空当前笔迹}只清除未保存的画布内容。已完成任务可通过\ui{继续／打开任务}重新打开纠错。

历史备份在采集目录的\code{symbol\_collection/revisions}，不参与字体制作。已有字体和冻结文档保持原样；修改采集样本后需要重新制作字体才能采用新笔迹。

类别未勾选时显示空心灰框，勾选后显示系统强调色底和白色对勾，便于在深色与浅色主题下识别。

符号目录按英文写法、数字、希腊字母及变体、括号书名号与定界符、标点连接符、单位量值和数学符号分组。类别行显示代表字符，展开可逐项选择。十三对明确列出的全角、半角同形符号各只录一次。可通过指定共用供体预览来源，并由用户逐项确认。

制作字体时自动为共用标点建立各自 Unicode 字符，笔迹不拉伸，全角形式保留全角前进宽度。旧样本可复用，旧字体与冻结文档不会自动改写；使用更新后的字体需在新版重新制作。完整字体同时选择汉字生成结果和符号样本，不要用仅符号的验收字体替代完整字体。

打开\ui{标点与扩展}，点击\ui{创建任务}。在弹窗中填写任务名称、作者 ID，勾选\ui{英文正体}、\ui{英文斜体}、数字、运算符、希腊字母、希腊变体、括号或标点等类别；设置每字符份数，默认五份，可设一至一百份。下方显示字符/写法数与总样本数；未选择类别时不能创建。英文正体和斜体可同时加入同一任务，同一字母的两种写法分别保存，继续任务后保持分类。取消配置不会替换当前任务。导入 TXT/CSV 时保留字表份数，TXT 与 CSV 空份数默认五份。画布大小可以调整，存在未保存笔迹时先保存或清空。符号采集结果通过字体库加入同一作者的字体，无需经过汉字生成。

正文优先选择正体，公式变量优先选择斜体；没有相应类别会提示，应用不会自动倾斜笔迹来伪造实录斜体。非标准笔数实录可以用于排版，不代表它具有可靠的标准笔画映射或部件生成资格。

\section{备份与设备验证}
保留原始采集 JSON 和任务记录，跨机器转移时确认作者 ID 与目录完整。生成结果和字体发布包分别保存，不覆盖原始样本。真实数位板压力、设备驱动及写字机纸面效果需要实际硬件验证；模拟采集检查不能代替这些验收。Windows 统一应用的构建与验收结果以对应版本记录为准。

\section{自由组合与确认复用（0.1.10）}
创建任务不提供快捷预设。各类字符默认不选，可整类勾选或展开逐字选择；英文大小写分别显示，斜体仅选需要补录的字母。默认每字五份，已有活动样本按同作者、字符和写法补差额；不会重建旧任务或清除进度。

点击\ui{查看笔迹并配置复用…}，查看每份已有英文正体笔迹，逐项确认斜体复用或相似希腊大写字符复用。未选录并不表示复用。未录来源须先采集；取消不保存，切换作者清除未提交关系。也可仅保存复用配置。制作字体须选择含 tasks 和 samples 的完整采集目录或 ZIP；来源缺失或重录时需重新确认。独立录制的目标字形优先，原始笔迹不变形，希腊字符保留自身身份。已有字体与冻结文档保持原样。

\section{公式数学轴信息（0.1.11）}
从原始样本制作字体时，数学轴参考信息保存在受哈希保护的核心来源摘要中，开发者转用户版和排版项目保存重开均保留。未审阅校准提示仍显示。旧字体如果缺少数学轴，需选择原始符号样本重新制作；无需重新录制。应用不会凭空填入默认轴值，也不会静默更新旧字体或冻结项目。
